\section{\texorpdfstring{$9\le n\le14$}{9<=n<=14} 的情形}
\label{sec:small-uniform}

本节把仅需全局整数数据即可处理的两个较大情形，与仍需精确有限验证的四个较小情形分开叙述。整个证明不读取历史搜索树，也不使用预存的轨道表。

对 $k\ge3$，以 $\ell_k$ 和 $c_k$ 分别表示极大 $k$ 点直线与极大
$k$ 点圆的数目，并令 $m_k=\ell_k+c_k$。每个点对属于唯一的极大直线；每个三点组
或者属于其极大直线，或者属于唯一的外接圆。因此
\begin{align}
 \ell_2+\sum_{k\ge3}\binom{k}{2}\ell_k&=\binom n2,
 \label{eq:small-line-pairs}\\
 \sum_{k\ge3}\binom{k}{3}m_k&=\binom n3,
 \qquad \sum_{k\ge3}c_k=c(P).
 \label{eq:small-triple-partition}
\end{align}
称 $(\ell_3,\ldots,\ell_K;c_3,\ldots,c_K)$ 为一个\emph{全局计数向量}。引理~\ref{lem:largest} 分别排除
$n=9,10,11,12,13,14$ 时大小至少为 $6,7,7,7,7,8$ 的块，故以下各分支只需考虑
$K\le5,6,6,6,6,7$。

计数向量的生成只使用必要条件。由于 $K<n-1$，删除任意一点后所得点集仍满足本问题的条件，对所有删除求和得到
\begin{equation}\label{eq:small-one-delete}
 n c(P)-3c_3\ge n c(n-1).
\end{equation}
在每一点反演后应用 Zhang 的指定外点普通线定理
\cite[Theorem~4.1]{zhang}，得到
\begin{equation}\label{eq:small-z-circle}
 3c_3\ge n\left\lceil\frac{n-1}{6}\right\rceil.
\end{equation}
此外，对计数向量的直线部分使用普通线、Melchior、Hirzebruch--Bojanowski 与 Shnurnikov 不等式，并在必定非共线的子集上求和
\cite[Eq.~(9)]{hirzebruch}\cite[Theorem~2.1]{pokora}\cite[Theorem~3]{shnurnikov}。
若子集包含在一条极大直线中，则不对它使用圆数下界；若子集包含在一个极大圆中，则只要求保留该圆。这是小阶删除平均中必须加入的修正。

还需用到块的交集所给出的下列初等不等式。对 $N\ge1$ 以及
$S=aN+b$（$0\le b<N$），令
\[
 \Phi_N(S)=N\binom a2+ab.
\]
它是满足 $\sum_i x_i=S$ 的非负整数向量中
$\sum_i\binom{x_i}{2}$ 的最小值。任取 $L$ 条极大直线和 $G$ 个极大圆，记它们的总点关联数和总点对关联数为 $I_1,I_2$，则
\begin{align}
 \Phi_n(I_1)&\le \binom L2+2LG+2\binom G2,
 \label{eq:small-block-point}\\
 \Phi_{\binom n2}(I_2)&\le LG+\binom G2.
 \label{eq:small-block-pair}
\end{align}
这是因为两条直线至多交于一点，而其余任意两个不同块至多交于两点。把每一类的
块按大小排列，并对每个初始段应用这两个不等式，即可覆盖所有子族；这一过程不必
决定哪些带标号的点实际落在所选直线或圆上。

\subsection{增广直线不等式}

固定 $x\in P$，以 $x$ 为中心反演，再对其余 $n-1$ 个像点作射影对偶。记所得直线排列中重数为 $j$ 的顶点数为 $t_j(x)$，并令
\[
 \delta_x=t_2(x)-3-\sum_{j\ge4}(j-3)t_j(x).
\]
以 $d_k^L(x)$ 表示经过 $x$ 的极大 $k$ 点直线数。

\begin{lemma}\label{lem:augmented-line}
对每个 $x\in P$，均有
\begin{equation}\label{eq:augmented-line-point}
 \sum_{k\ge3}k d_k^L(x)\le n-1+\delta_x.
\end{equation}
因此
\begin{equation}\label{eq:augmented-line-summed}
 \sum_{k\ge3}k^2\ell_k\le n(n-1)+D,
 \qquad
 D=3m_3-3n-\sum_{k\ge5}k(k-4)m_k.
\end{equation}
\end{lemma}

\begin{proof}
记上述 $n-1$ 条对偶直线的排列为 $\mathcal A_x$，再加入反演中心的对偶直线。经过 $x$ 的一条原 $k$ 点直线，对应于新增直线上的一个旧的 $(k-1)$ 重顶点。设这些旧顶点的重数为 $i_1,\ldots,i_s$。新增直线产生
$(n-1)-\sum_j i_j$ 个新的普通顶点；每把一个旧顶点的重数提高一，Melchior 亏量减少一，这对二重和三重顶点也成立。因此增广排列的亏量为
\[
 \delta_x+(n-1)-\sum_{j=1}^s(i_j+1)
 =\delta_x+(n-1)-\sum_{k\ge3}k d_k^L(x).
\]
增广排列是本质排列，故 Melchior 不等式给出
\eqref{eq:augmented-line-point}。再对 $x$ 求和，并使用
$\sum_xd_k^L(x)=k\ell_k$ 及局部关联恒等式，即得
\eqref{eq:augmented-line-summed}。
\end{proof}

\subsection{十三点与十四点的全局排除}

当 $n=14$ 时，在 $c(P)<F(14)$ 下，前述必要条件的全部整数解共有 8981 个。
不等式~\eqref{eq:augmented-line-summed} 排除其中 8980 个；唯一剩余的计数向量为
\[
 (\ell_3,\ldots,\ell_7;c_3,\ldots,c_7)
 =(20,0,0,0,0;23,44,0,2,3).
\]
其中三个七点圆两两至多交于两点，故其并至少含
$3\cdot7-3\cdot2=15$ 点，矛盾。因此 $c(14)=F(14)$。

当 $n=13$ 时，小于 $F(13)$ 的整数计数向量共有 3962 个。增广直线不等式后剩 137 个，再由
\eqref{eq:small-block-point}--\eqref{eq:small-block-pair} 剩 73 个。以下两条不等式排除全部剩余向量。

在任一点处，局部排列含十二条直线；在当前分支中，其顶点重数不超过五。将其重数向量记为
$(t_2,t_3,t_4,t_5)$，亏量记为 $\delta$。则
\begin{equation}\label{eq:n13-local-two-ineq}
 -2\delta+t_4+3t_5\le6,
 \qquad -\delta+t_5\le1.
\end{equation}
为证明第一式，由点对恒等式及 Bojanowski 不等式有
\begin{align*}
 t_2+3t_3+6t_4+10t_5&=66,\\
 4t_2+3t_3&\ge48+5t_5,
\end{align*}
从而
\[
 H:=3t_2-6t_4-15t_5+18\ge0.
\]
令 $E=2t_2-3t_4-7t_5$。由 $H+3\delta=3(E+3)$ 得
$E\ge-3$；点对恒等式又给出 $E\equiv0\pmod3$。若 $E=-3$，则必有
$H=\delta=0$。Cuntz 对单纯十二线排列的完备分类在排除六重顶点后只留下
$(8,10,3,1)$ 和 $(9,7,6,0)$，而二者均满足 $E=0$
\cite[Theorem~1.1 and the invariant tables]{cuntz}。故 $E\ge0$，这等价于
\eqref{eq:n13-local-two-ineq} 的第一式。

为证明第二式，四个五重顶点对偶为四条五点线，其并至少含
$20-\binom42=14$ 点，故 $t_5\le3$。若 $\delta\ge2$，结论立即成立；若
$\delta=0$，上述 Cuntz 类型均满足 $t_5\le1$。唯一尚可能违反不等式的是
$(\delta,t_5)=(1,3)$。此时点对恒等式与亏量恒等式迫使
$(t_2,t_3,t_4,t_5)=(12,4,2,3)$。但三条五点线与一条四点线的并至少含
$5+5+5+4-\binom42=13$ 点，再次矛盾。第二式得证。

对十三个中心求和，\eqref{eq:n13-local-two-ineq} 化为
\begin{equation}\label{eq:n13-global-two-ineq}
 -2D+5m_5+18m_6\le78,
 \qquad -D+6m_6\le13.
\end{equation}
在最后 73 个计数向量中，71 个违反第一式，其余两个违反第二式。这里仅是把整数计数向量代入
\eqref{eq:n13-global-two-ineq}，没有选择任何带标号点子集。因此 $c(13)=F(13)$。

\subsection{九点至十二点的精确验证}

对 $9\le n\le12$，同一组全局条件与
\eqref{eq:augmented-line-summed} 留下少量非空边界。下面给出统一验证的每一步。固定一点 $x$，以 $a_k(x),u_k(x)$ 分别表示经过 $x$ 的极大
$k$ 点直线和圆的数目，则
\[
 t_{k-1}(x)=a_k(x)+u_k(x),
 \qquad
 \sum_{j\ge2}\binom j2t_j(x)=\binom{n-1}{2}.
\]
这些局部重数向量满足前述公开的直线排列不等式、块交集不等式及
\eqref{eq:augmented-line-point}。第一个整数模型只记录每种局部重数向量出现多少次，并要求
\[
 \sum_xa_k(x)=k\ell_k,\qquad
 \sum_xu_k(x)=kc_k,\qquad
 \sum_x\delta_x=D.
\]

后两个阶段仍不指定带标号的点子集。对块类别 $\alpha$，令 $d_\alpha(x)$ 为包含
点 $x$ 的该类块数，令
$\lambda_\alpha(e)$ 为点对 $e$ 上的度；若块大小为 $k_\alpha$、数目为
$N_\alpha$，则
\begin{align}
 \sum_xd_\alpha(x)&=k_\alpha N_\alpha,&
 \sum_e\lambda_\alpha(e)&=\binom{k_\alpha}{2}N_\alpha,
 \label{eq:small-category-totals}\\
 \sum_\alpha(k_\alpha-2)\lambda_\alpha(e)&=n-2.
 \label{eq:small-pair-profile2}
\end{align}
最后一个恒等式来自：经过 $e$ 的所有极大块在 $e$ 外的部分恰好分割
$P\setminus e$。点数据与点对数据还须满足
\[
\sum_{y\ne x}\lambda_\alpha(\{x,y\})
=(k_\alpha-1)d_\alpha(x).
\]
随后，对任意两个块类别使用精确的交集计数恒等式
\[
 N_2=M_2,\qquad N_1=M_1-2M_2,\qquad N_0=Q-M_1+M_2,
\]
其中 $M_1$ 是所有不同块对交集大小之和，$M_2$ 是这些交集大小的二阶二项式
系数之和，$N_s$ 是恰交于 $s$ 点的不同块对数；模型还精确计数与一个指定块不交的
三点组。所有系数均来自二项式计数，不使用浮点松弛。

最后，以 $x$ 为中心反演会确定一个完整的 $(n-1)$ 点像集计数向量：经过
$x$ 的块变为少一个点的像直线，不经过 $x$ 的块变为像圆。只有当该子向量在下一层
未被排除时，当前局部重数向量才被保留；随后再次平衡全部 $n$ 个点处的数据。递归在
$n=9$ 终止。九点模型在带标号点集的所有子集中选择极大直线与圆，要求每个三点组恰属于一个块、每个点对至多属于一条所选直线，并对每个满足本问题条件的真子集使用已经证明的圆数下界。这个抽象模型比几何实现更宽松，故其不可行足以推出几何不可行。

最终重放的精确数量如下：
\begin{center}
\small
\begin{tabular}{c|rrrrrrr}
\toprule
$n$ & 初始 & 增广求和 & 局部向量 & 点--点对 & 交集计数 & 块对 & 终点\\
\midrule
9  &45   &36  &23  &-- &-- &-- &0\\
10 &180  &150 &55  &26 &21 &12 &1\\
11 &416  &160 &70  &39 &28 &-- &0\\
12 &1354 &211 &169 &55 &44 &-- &0\\
\bottomrule
\end{tabular}
\end{center}
“初始”一列分别由 129、551、1215、4322 个原始计数向量缩减而来。
$n=11,12$ 的终点零由同一个递归器得到；它们分别访问
\[
 (28,72,222)\quad\text{与}\quad(44,322,661,622)
\]
个 $(11,10,9)$ 阶及 $(12,11,10,9)$ 阶计数向量。每个被排除的状态均为精确的
\texttt{INFEASIBLE}，没有 \texttt{UNKNOWN}。

当 $n=10$ 时，递归只留下一个完整计数向量
\[
 (\ell_3,\ell_4,\ell_5,\ell_6;
   c_3,c_4,c_5,c_6;\ell_2)
 =(10,0,0,0;10,20,2,0;15).
\]
递归证书保留的四种局部重数向量均满足 $u_5(x)=1$ 与 $a_3(x)\le3$。两个五点圆共有十个
点关联，故它们恰好把十点分割为 $A,B$。十条三点线共有三十个点关联，故每点均有
$a_3(x)=3$；相应的唯一平衡局部向量还满足 $u_4(x)=8$。每个四点圆因而在 $A,B$ 两侧
各取两点。

把这样的四点圆对应于其在 $A$ 与 $B$ 中的两条弦。两弦的交点位于两个五点圆的
固定根轴上。圆上五点所确定的弦中，不可能有三条在圆外同点：三条两两不共端点的
弦需要六个端点。因此，在根轴同一点相交的两侧弦数都至多为二。二十个四点圆达到
该上界，弦图遂恰为五个互不相交的 $K_{2,2}$。在 $S_5\times S_5$ 作用下，
Petersen 图的六个完美匹配给出两个双陪集轨道。其中一个不存在相容的十条三点线
设计，另一个有两个互补设计，二者属于同一轨道。

对后一个轨道取代表元，十条直线关联给出含非零参数 $a,b,d$ 的完备射影参数化。选取五个圆交点矩阵的三阶子式，得到
\begin{gather*}
 ad+b^2,\quad -ab+ad-b,\quad ab-ad-a-1,\\
 -ad+bd-b-d,\quad b(ab+1).
\end{gather*}
它们必须同时为零。第二、三式相加给出 $a+b+1=0$，而末式及 $b\ne0$ 给出
$ab+1=0$，故 $b^2+b-1=0$。第一式给出
$d=b^2/(b+1)$；这里 $b+1=-a\ne0$。再代入第四式并移去非零因子，得到
\[
 (b-1)(2b+1)=0,
\]
它与 $b^2+b-1$ 互素。被舍弃的参数分支会使两个带标号点重合；形成原始子式时除去的因子只可能是
$a,d,b,ad$，它们均非零。因此最后一个十点计数向量也不可能实现。

结合式~\eqref{eq:F} 的构造，得到
\[
 c(n)=F(n)\qquad(9\le n\le14).
\]

\begin{remark}[有限验证的范围]
本节只有九点终止模型实际选择带标号点的子集。十点至十二点只使用计数向量及
局部重数的递归一致性，十三点与十四点只使用全局不等式。证明不对任意直线排列假设
$\sum_jt_j\ge2r-4$；对含 $q$ 条直线的排列，Elliott 的 Theorem~3 仅在明确满足
$q\ge10$ 且不存在 $(q-1)$ 重顶点时使用。
\end{remark}
